{\footnotesize
\setlength{\tabcolsep}{3.5pt}
\begin{threeparttable}
\begin{tabular}{lllrrrrr}
\toprule
 & Region & State & Initial employment & $T$ & $D$ & $F$ & Portab. \\
\midrule
Most flexible & Rio de Janeiro & Rio de Janeiro & 2,973,469 & 0.078 & 0.028 & 2.81 & 0.680 \\
 & Porto Alegre & Rio Grande do Sul & 1,264,900 & 0.049 & 0.018 & 2.77 & 0.655 \\
 & Curitiba & Paraná & 1,112,706 & 0.056 & 0.022 & 2.55 & 0.706 \\
 & Campinas & São Paulo & 1,920,219 & 0.050 & 0.020 & 2.54 & 0.701 \\
 & São Paulo & São Paulo & 6,928,315 & 0.062 & 0.025 & 2.47 & 0.707 \\
 & Altamira & Pará & 53,603 & 0.491 & 0.203 & 2.42 & 0.706 \\
 & Caxias do Sul & Rio Grande do Sul & 393,565 & 0.070 & 0.030 & 2.31 & 0.694 \\
 & Oiapoque - Porto Grande & Amapá & 6,961 & 0.247 & 0.115 & 2.15 & 0.720 \\
 & São José dos Campos & São Paulo & 527,454 & 0.073 & 0.034 & 2.14 & 0.678 \\
 & Belo Horizonte & Minas Gerais & 1,673,843 & 0.061 & 0.029 & 2.11 & 0.699 \\
\addlinespace[3pt]
Least flexible & Redenção & Pará & 39,702 & 0.113 & 0.717 & 0.16 & 0.576 \\
 & Sousa - Cajazeiras & Paraíba & 14,850 & 0.133 & 0.942 & 0.14 & 0.647 \\
 & Rio Branco & Acre & 67,856 & 0.106 & 1.050 & 0.10 & 0.654 \\
 & Barreiras & Bahia & 65,399 & 0.080 & 0.827 & 0.10 & 0.639 \\
 & São Raimundo Nonato & Piauí & 4,530 & 0.291 & 4.713 & 0.06 & 0.609 \\
 & Caruaru & Pernambuco & 163,639 & 0.108 & 1.847 & 0.06 & 0.632 \\
 & Quixadá & Ceará & 52,645 & 0.189 & 3.525 & 0.05 & 0.509 \\
 & Crateús & Ceará & 7,214 & 0.179 & 3.395 & 0.05 & 0.603 \\
 & Manaus & Amazonas & 421,620 & 0.096 & 2.409 & 0.04 & 0.699 \\
 & Corrente - Bom Jesus & Piauí & 6,101 & 0.097 & 5.432 & 0.02 & 0.602 \\
\bottomrule
\end{tabular}
\begin{tablenotes}[flushleft]\footnotesize
\item \textit{Note:} Intermediate regions ranked by allocative flexibility in the 2014 to 2019 window. Source: RAIS, authors' calculations.
\end{tablenotes}
\end{threeparttable}
}
